\section{Power energy controls all original degrees and unbounded energies}
\label{sec:r11-power-all-degree}

Keep the explicit domain condition (DG1): the compact closure of
the ordinary open set $V$ lies in the ordinary interior of one
original $K$-affinoid domain. The background $S$ is an arbitrary
original symmetric middle-positive separately closed $(p,p)$-current
on $V$, and $q=n-p\geq1$. No smoothness, geometric origin or prior
BT representation of $S$ is assumed.

Fix $\sigma>0$ and a finite family
$u_1,\ldots,u_m\in\mathcal E^\sigma_{S,\partial}(V)$. Put
\[
 \kappa=q+\sigma,\quad w=\sum_i u_i,\quad W=1-w,\quad
 W_j=1-\sum_i\max(u_i,-j),\quad E_j=\{w>-j\}.
 \tag{PG1}
\]
Let $T_{\alpha,j}$ and $N_\alpha$ be the bounded and unbounded
mixed currents from Section~\ref{sec:r11-capacity-dimension-gain};
$|\alpha|=k\leq q$ includes repeated slots and $\alpha=0$.
The power-energy cone theorem and (PE2) give
\[
 {\cal J}_w(R):=\int_R^\infty
       t^{\kappa-1}C_S(\{w<-t\};V)\,dt<\infty
       \quad(R>0),\qquad {\cal J}_w(R)\longrightarrow0.
 \tag{PG2}
\]
This is an actual capacity integral in the fixed domain $V$.

\begin{lemma}[Graded coefficient moments including the endpoint]
\label{lem:r11-graded-curvature-moments}
Let $a\geq0$ and $a+k\leq\kappa$. For every original compact
positive coefficient control $\Theta$ of degree $(q-k,q-k)$,
the measures $W^a\Theta\wedge N_\alpha$ are locally finite.
There is a finite constant $C_\Theta$, independent of $j,R$, such
that for $R\geq1$,
\[
 \begin{split}
 \sup_{j\geq1}\int_{\{w\leq-R\}}W_j^a\,
                           \Theta\wedge T_{\alpha,j}
     &\leq C_\Theta R^{a+k-\kappa}{\cal J}_w(R/4),\\
 \int_{\{w\leq-R\}}W^a\,\Theta\wedge N_\alpha
     &\leq C_\Theta R^{a+k-\kappa}{\cal J}_w(R/4).
 \end{split}
 \tag{PG3}
\]
The constant may depend on the family size, degree, exponent,
and coefficient control. Each full local moment in the first
line, with the tail indicator removed, is uniformly bounded in $j$.

For every continuous $F$ on the nonpositive orthant with
$|F(x)|\leq A(1-\sum_i x_i)^a$, there is local coefficient
total-variation convergence
\[
               F(\mathbf u_j)T_{\alpha,j}
                          \longrightarrow F(\mathbf u)N_\alpha.
 \tag{PG4}
\]
The endpoint $a+k=\kappa$ is included for this fixed family.
Multipliers on its pole union are immaterial because all these
currents have zero coefficient mass there.
\end{lemma}

\noindent\textbf{Status: \texttt{NEW\_PROVED}} as a
degree-dependent strengthening of (DG9)--(DG11), using the
proved power-energy capacity criterion.

\begin{proof}
Write $F_w(t)=C_S(\{w<-t\};V)$. On
$A_T=\{-2T<w\leq-T\}$ all $u_i>-2T$. For any $j$, strict
bounded maximum locality permits replacing each $u_{i,j}$
on this shell by $\max(u_i,-\min(j,2T))$, whose global range
is $[-2T,0]$. Estimate (DG3) and $W_j\leq W\leq1+2T$ give
\[
 \int_{A_T}W_j^a\Theta\wedge T_{\alpha,j}
       \leq C T^{a+k}F_w(T/2).
 \tag{PG5}
\]
This is uniform even when $j<T$. Each finite-stage current is
zero on the pole union by (DG3). The same shell estimate holds
for $N_\alpha$ by consistent restriction. Monotonicity gives
\[
 T^\kappa F_w(T/2)
       \leq C_\kappa\int_{T/4}^{T/2}
                         t^{\kappa-1}F_w(t)\,dt.
 \tag{PG6}
\]
Sum (PG5) over $T=2^\ell R$, $\ell\geq0$. The factors
$T^{a+k-\kappa}$ are bounded by $R^{a+k-\kappa}$, and the
integration intervals in (PG6) are disjoint. This proves (PG3).
On $\{w>-1\}$ consistency with the fixed level-one product
controls the remaining compact coefficient masses.

The bounded and limiting weighted currents agree on $E_j$.
Their variations on its complement are bounded by (PG3) at
$R=j$ and the growth bound for $F$. The bounds tend to zero,
also at equality, because $ {\cal J}_w(j/4)\to0$.
Finite coefficient domination proves (PG4) for every original
test and actual refinement.
\end{proof}

\begin{theorem}[Full recursive calculus in every positive power class]
\label{thm:r11-power-full-recursion}
For every $\sigma>0$, every finite family in
$\mathcal E^\sigma_{S,\partial}(V)$, and every
$|\alpha|\leq q-1$, the multiplier $u_iN_\alpha$ is defined
on all original coefficients and
\[
                         D(u_iN_\alpha)=N_{\alpha+e_i}.
 \tag{PG7}
\]
All repeated and mixed products in this class therefore have
the full recursive differential meaning through top degree,
in addition to separate closedness, symmetry, multilinearity,
and truncation independence.
\end{theorem}

\noindent\textbf{Status: \texttt{NEW\_PROVED}} as an
all-degree consequence of one top power-energy condition per
potential. The affinoid domain hypothesis remains in force.

\begin{proof}
For $k=|\alpha|\leq q-1$, one has $1+k\leq q<q+\sigma$.
Apply (PG4) with $F(x)=x_i$ before differentiating. The bounded
identity $D(u_{i,j}T_{\alpha,j})=T_{\alpha+e_i,j}$ then passes
distributionally to (PG7). The other properties were proved
without this recursion in the dimension-gain theorem; there
is no circular definition of the top product.
\end{proof}

\subsection{A direct capacity bound for bounded energy currents}

\begin{lemma}[The extra two powers in an energy trace]
\label{lem:r11-bounded-energy-capacity-power}
Let $-L\leq b_i\leq0$ be bounded classical psh functions on $V$,
and let $Q_\alpha$ be their bounded mixed product over $S$,
$|\alpha|=k\leq q-1$. For every original compact positive
control $\Theta$ of degree $(q-k-1,q-k-1)$ and Borel $A\subset V$,
\[
       \int_A\Theta\wedge\Gamma_{Q_\alpha}(b_i)
                 \leq C_\Theta L^{k+2}C_S(A;V).
 \tag{PG8}
\]
The constant is independent of the bounded potentials and $L>0$.
Mixed energy coefficient variations have the corresponding
bound by the sum of the diagonal traces.
\end{lemma}

\noindent\textbf{Status: \texttt{DIRECT\_COROLLARY}} of
the bounded energy identity and the global original-coefficient
compensator, retaining the arbitrary background.

\begin{proof}
Put $v=b_i+L$, so $0\leq v\leq L$ and $v^2$ is bounded
classical psh. The energy identity gives
\[
 0\leq\Gamma_{Q_\alpha}(b_i)
       =\tfrac12D(v^2Q_\alpha)-vD(b_iQ_\alpha)
       \leq\tfrac12D(v^2Q_\alpha).
 \tag{PG9}
\]
These are middle-cone inequalities; the nonnegative Borel
multiplier in the second term preserves positivity.
Use (DG2) for the energy current and its original test.
Normalize the $k$ slots in $Q_\alpha$ by $L$, the first
slot $v^2$ by $L^2$, and fill the remaining $q-k-1$ slots
with the global normalized smooth compensator. Positive
mixed expansion gives the factor $(q^q/q!)C_S$ and proves
(PG8). The positive energy matrix and its measure
Cauchy--Schwarz inequality control cross variations.
No energy current has been assumed closed.
\end{proof}

\begin{theorem}[Unbounded energy currents in the graded power domain]
\label{thm:r11-power-graded-energy}
Assume $|\alpha|=k\leq q-1$ and $k+2\leq\kappa=q+\sigma$.
There is an intrinsic matrix of coefficient-Radon currents
\[
 G^\alpha_{ia}
  =\tfrac12\left[
     D(u_i u_aN_\alpha)-u_iN_{\alpha+e_a}
                           -u_aN_{\alpha+e_i}\right].
 \tag{PG10}
\]
It is positive semidefinite in the original middle cone,
has zero coefficient mass on the pole union, and satisfies
\[
       \Gamma_{T_{\alpha,j}}(u_{i,j},u_{a,j})
                                \longrightarrow G^\alpha_{ia}
 \tag{PG11}
\]
in local coefficient total variation. If $b\geq0$ and
$b+k+2\leq\kappa$, then for every original compact positive
control $\Theta$,
\[
 \begin{split}
 \sup_j\int_{\{w\leq-R\}} W_j^b
           \Theta\wedge\Gamma_{T_{\alpha,j}}(u_{i,j})
   &\leq C_\Theta R^{b+k+2-\kappa}{\cal J}_w(R/4),\\
 \int_{\{w\leq-R\}} W^b\Theta\wedge G^\alpha_{ii}
   &\leq C_\Theta R^{b+k+2-\kappa}{\cal J}_w(R/4).
 \end{split}
 \tag{PG12}
\]
The corresponding full local moments are finite. A continuous
weight of growth at most $W^b$ may be inserted in (PG11),
with the same total-variation conclusion for this fixed tuple.

In particular $\sigma\geq1$ gives these energy currents in
every permitted energy degree. If $0<\sigma<1$, the conclusion
applies to $k\leq q-2$. The energy currents need not be closed.
\end{theorem}

\noindent\textbf{Status: \texttt{NEW\_PROVED}}. This improves
the sufficient superquadratic all-trace hypothesis (UE2) on the
printed domain, and includes the finite first-power energy endpoint.

\begin{proof}
On $A_T=\{-2T<w\leq-T\}$ replace every truncated slot by its
further truncation at $-2T$. Lemma
\ref{lem:r11-energy-truncation-locality} permits this in the
background slots and the differentiated energy slot. Its proof
uses shifted squares and four-branch maxima, so also applies
with unequal initial truncation levels. Estimate (PG8) gives
\[
 \int_{A_T}W_j^b
        \Theta\wedge\Gamma_{T_{\alpha,j}}(u_{i,j})
                  \leq C T^{b+k+2}F_w(T/2).
 \tag{PG13}
\]
Each finite-stage energy is zero on the pole union by (PG8)
and $C_S(P;V)=0$. Sum (PG13) with (PG6) to obtain the first
bound of (PG12).

Bounded energy locality gives consistency of the restrictions
on $E_j$. The positive cut diagonal energies increase, and
the shell estimates and level-one core give locally finite
limits. Their tails and those of the uncut energies satisfy
(PG12), first with $b=0$, then with the stated weights.
Both tails tend to zero, proving total-variation convergence.
Polarization gives the cross limits, the positive matrix
assertion, and all coefficient variation bounds.

The three terms in (PG10) are defined before identification:
use (PG3) with $(a,k)=(2,k)$ and $(a,k)=(1,k+1)$.
Their total exponent is $k+2\leq\kappa$. Pass the bounded
three-term energy formula through (PG4) and distributional
continuity. This identifies the limits with (PG10), including
every original compact test and actual refinement.
No Borel-restricted current has been differentiated.
\end{proof}

\begin{theorem}[A degree-dependent unbounded chain rule]
\label{thm:r11-power-graded-chain}
Assume the preceding hypotheses. Let $F$ be $C^2$ near
the nonpositive orthant and suppose, for some
$2\leq r\leq\kappa-k$ and finite $A$,
\[
 |F(x)|\leq A W(x)^r,\quad
 |F_i(x)|\leq A W(x)^{r-1},\quad
 |F_{ia}(x)|\leq A W(x)^{r-2},
 \qquad W(x)=1-\sum_i x_i.
 \tag{PG14}
\]
Then the following is an equality of original
coefficient-defined currents on ordinary opens:
\[
 D(F(\mathbf u)N_\alpha)
       =\sum_i F_i(\mathbf u)N_{\alpha+e_i}
                         +\sum_{i,a}F_{ia}(\mathbf u)G^\alpha_{ia}.
 \tag{PG15}
\]
Convex coordinatewise nondecreasing $F$ gives a middle-positive
right side. The endpoint $r=\kappa-k$ is permitted for each
fixed finite-power-energy family.
\end{theorem}

\noindent\textbf{Status: \texttt{NEW\_PROVED}} with the
printed growth, background and domain hypotheses.

\begin{proof}
Apply the bounded multivariable chain identity to common
truncations. The term to be differentiated has growth $W^r$
and degree $k$; each curvature multiplier has growth
$W^{r-1}$ and degree $k+1$; each energy multiplier has growth
$W^{r-2}$ and energy budget $k+2$. Each total exponent is
$r+k\leq\kappa$. Equations (PG4) and (PG12) therefore pass
every term in local coefficient total variation.
Distributional continuity proves (PG15). The sign assertion
uses $F_i\geq0$, the positive Hessian of $F$, and the positive
energy matrix, without specializing $S$.
\end{proof}

\paragraph{Uniform families and the critical boundary.}
For a uniform $\sigma$-power energy bound, (PE10) and (PE4)
bound $ {\cal J}_{w_l}(0)$ uniformly. The tail estimates above
are then uniform and vanish whenever the relevant total
exponent is strictly below $q+\sigma$. At equality, a bound
for the full integral does not force a uniformly vanishing
tail. This distinction is essential for varying tuples;
it does not affect the fixed-tuple endpoint constructions.
